%% ma: L10-B13-0080
%% lop: 10
%% chuong: 5
%% bai: 13
%% dang: 10.13.10
%% loai: TN4
%% muc: TH
%% dap_an: "A"
%% nguon: "chuyen-de-cac-so-dac-trung-cua-mau-so-lieu-khong-ghep-nhom-toan-10-chuong-trinh-moi.pdf (GV. Phan Nhật Linh) (P03), Bài 02 Dạng 2 Phần I Câu 5"
%% kien_thuc: "Lớp 10 Bài 13 (trung vị, tứ phân vị)"
%% trang_thai: chinh-thuc
%% ghi_chu: "Tự duyệt theo yêu cầu giáo viên 22/10/2026."
\begin{ex}
Bảng số liệu sau cho biết sản lượng chè thu được trong $1$ năm (kg/sào) của $16$ hộ gia đình:
\begin{center}
\begin{tabular}{|c|c|c|c|c|c|c|c|}\hline
$111$ & $112$ & $113$ & $112$ & $114$ & $127$ & $128$ & $125$\\\hline
$119$ & $118$ & $113$ & $126$ & $120$ & $115$ & $123$ & $116$\\\hline
\end{tabular}
\end{center}
Các tứ phân vị của mẫu số liệu là
\choice{\True $Q_1=113;Q_2=117;Q_3=124$}{$Q_1=117;Q_2=113;Q_3=124$}{$Q_1=113;Q_2=117;Q_3=123$}{$Q_1=113;Q_2=122;Q_3=123$}
\loigiai{Sắp xếp: $111$; $112$; $112$; $113$; $113$; $114$; $115$; $116$; $118$; $119$; $120$; $123$; $125$; $126$; $127$; $128$.

$Q_2=\dfrac{116+118}2=117$; $Q_1=\dfrac{113+113}2=113$; $Q_3=\dfrac{123+125}2=124$. Chọn A.}
\end{ex}
